\def\uadectitulo{Pensamiento matemático}
\documentclass[12pt,letterpaper,oneside]{article}
\usepackage[T1]{fontenc}
\usepackage[utf8]{inputenc}
\usepackage[spanish,es-nodecimaldot]{babel}
\usepackage{lmodern}
\usepackage[margin=2.5cm]{geometry}
\usepackage{amsmath,amssymb,graphicx,booktabs,longtable,array,enumitem}
\usepackage[most]{tcolorbox}
\usepackage{tikz}
\usetikzlibrary{babel,arrows.meta,positioning}
\usepackage{setspace}
\usepackage[hidelinks]{hyperref}
\onehalfspacing
\setlength{\emergencystretch}{3em}
\providecommand{\uadectitulo}{Material de estudio}
\providecommand{\uadecautor}{Material de estudio AulaTeX}
\newcommand{\portadaUAdeC}{%
  \hypersetup{pageanchor=false}
  \begin{titlepage}
  \centering
  {\Large Universidad Autónoma de Coahuila\par}
  \vspace{0.7cm}
  {\large Preparación para ingreso a Médico Cirujano\par}
  \vspace{1.8cm}
  {\LARGE\bfseries\uadectitulo\par}
  \vspace{1cm}
  {\large Actividad 1: síntesis y representación conceptual\par}
  {\large Actividad 2: práctica y soluciones justificadas\par}
  \vfill
  {\uadecautor\par}
  {Revisión: 10 de octubre de 2026\par}
  \vspace{0.8cm}
  \begin{minipage}{0.87\linewidth}
  \small Material formativo local. No constituye una guía, plantilla o examen oficial de UAdeC o CENEVAL. La cobertura y dificultad no están certificadas frente a una convocatoria de ingreso. No sustituye orientación ni atención clínica.
  \end{minipage}
  \end{titlepage}
  \hypersetup{pageanchor=true}
  \tableofcontents\clearpage
  \section*{Alcance y procedencia}
  Las actividades siguen las consignas locales conservadas en esta materia. Los reactivos son de práctica y las respuestas son propuestas razonadas, no resultados de un intento del aspirante. Las fuentes de apoyo se identifican al final. Los PDF bibliográficos locales se inventarían por archivo y hash; su disponibilidad no acredita lectura íntegra ni permiso de redistribución.
}
\begin{document}
\portadaUAdeC
\section{Actividad 1: síntesis de investigación}
\subsection{Proporcionalidad y unidades}
Una razón compara cantidades. En proporcionalidad directa el cociente permanece y en inversa permanece el producto bajo las condiciones del modelo. Un porcentaje usa base cien; aumentos sucesivos se multiplican, no se suman sin revisar la base.
\subsection{Álgebra y geometría}
Resolver una ecuación implica operaciones equivalentes a ambos lados y verificar el resultado. Área y perímetro tienen unidades distintas. Para triángulos rectángulos se aplica Pitágoras; no debe aplicarse sin esa condición.
\subsection{Datos y probabilidad}
Media suma y divide por cantidad; mediana requiere ordenar. La probabilidad clásica necesita resultados equiprobables. Los ejemplos siguientes son ficticios y no describen mediciones clínicas.
\subsection{Representación conceptual y cuadro de síntesis}
\begin{longtable}{p{0.27\linewidth}p{0.65\linewidth}}
\toprule Concepto & Relación o función\\\midrule
Datos & Cantidad, unidad y aulaTEX/base\\
Modelo & Razón, ecuación o distribución\\
Cálculo & Operaciones justificadas\\
Comprobación & Sustitución, unidades y plausibilidad\\
\bottomrule
\end{longtable}
\begin{figure}[htbp]
\centering
\begin{tikzpicture}[node distance=5mm, every node/.style={draw,rounded corners=2pt,text width=10cm,align=center,font=\small}]
\node (step0) {Datos: Cantidad, unidad y base};
\node[below=of step0] (step1) {Modelo: Razón, ecuación o distribución};
\draw[-{Stealth}] (step0) -- (step1);
\node[below=of step1] (step2) {Cálculo: Operaciones justificadas};
\draw[-{Stealth}] (step1) -- (step2);
\node[below=of step2] (step3) {Comprobación: Sustitución, unidades y plausibilidad};
\draw[-{Stealth}] (step2) -- (step3);
\end{tikzpicture}
\caption{Esquema didáctico original; las flechas indican orden de lectura o relación conceptual, no causalidad universal.}
\end{figure}
\clearpage\section{Actividad 2: cuestionario original resuelto}
\begin{enumerate}
\item\textbf{Tres cuadernos cuestan 45 pesos. Al mismo precio unitario, ¿cuánto cuestan ocho?}
\begin{enumerate}[label=\Alph*)]
\item 80
\item 100
\item 120
\item 135
\end{enumerate}
\textbf{Respuesta: C.} 45/3 = 15 pesos por cuaderno; 8 por 15 = 120. Las otras cantidades no conservan el precio unitario.
\item\textbf{Un artículo de 200 pesos aumenta 10\% y luego baja 10\% sobre el nuevo precio. ¿Precio final?}
\begin{enumerate}[label=\Alph*)]
\item 198
\item 200
\item 180
\item 220
\end{enumerate}
\textbf{Respuesta: A.} 200 por 1.10 por 0.90 = 198. Los porcentajes usan bases diferentes y no se cancelan.
\item\textbf{Resuelve 3x + 5 = 20.}
\begin{enumerate}[label=\Alph*)]
\item 3
\item 4
\item 6
\item 5
\end{enumerate}
\textbf{Respuesta: D.} Restar 5 y dividir entre 3 da x = 5; sustitución: 15 + 5 = 20.
\item\textbf{¿Mediana de 2, 4, 9, 10 y 15?}
\begin{enumerate}[label=\Alph*)]
\item 8
\item 9
\item 10
\item 15
\end{enumerate}
\textbf{Respuesta: B.} Ordenados, el tercer valor es 9. Media y máximo no son mediana.
\item\textbf{En una bolsa hay tres fichas rojas y dos azules equiprobables. ¿Probabilidad de extraer una roja?}
\begin{enumerate}[label=\Alph*)]
\item 2/5
\item 1/2
\item 3/5
\item 3/2
\end{enumerate}
\textbf{Respuesta: C.} Tres favorables entre cinco posibles: 3/5; 3/2 excede uno.
\item\textbf{¿Área de un rectángulo de 6 cm por 4 cm?}
\begin{enumerate}[label=\Alph*)]
\item 24 cm cuadrados
\item 20 cm
\item 10 cm cuadrados
\item 24 cm
\end{enumerate}
\textbf{Respuesta: A.} Área = 6 por 4 = 24 cm cuadrados. 20 cm sería perímetro; cm solos no son unidad de área.
\item\textbf{¿Hipotenusa de un triángulo rectángulo con catetos 3 y 4 cm?}
\begin{enumerate}[label=\Alph*)]
\item 7 cm
\item 12 cm
\item 25 cm
\item 5 cm
\end{enumerate}
\textbf{Respuesta: D.} Raíz de 3 al cuadrado más 4 al cuadrado = raíz de 25 = 5.
\item\textbf{¿Media de 4, 6 y 8?}
\begin{enumerate}[label=\Alph*)]
\item 5
\item 6
\item 7
\item 18
\end{enumerate}
\textbf{Respuesta: B.} Suma 18 dividida entre tres = 6; 18 es la suma, no la media.
\item\textbf{Una tarea requiere 24 persona-horas. Si cuatro personas rinden igual simultáneamente, ¿tiempo ideal?}
\begin{enumerate}[label=\Alph*)]
\item 96 horas
\item 24 horas
\item 6 horas
\item 4 horas
\end{enumerate}
\textbf{Respuesta: C.} 24/4 = 6 horas bajo el modelo ideal indicado; no es garantía de productividad real.
\item\textbf{Una cantidad pasa de 50 a 65. ¿Aumento porcentual respecto a 50?}
\begin{enumerate}[label=\Alph*)]
\item 30\%
\item 15\%
\item 23\%
\item 65\%
\end{enumerate}
\textbf{Respuesta: A.} (65 - 50)/50 por 100 = 30\%; 15 es aumento absoluto, no porcentaje.
\end{enumerate}

\section{Conclusiones}
El material permite repasar conceptos y contrastar respuestas. La clave no acredita desempeño del aspirante ni aprobación institucional. Deben revisarse la convocatoria y el temario aplicables antes de usarlo como plan definitivo de ingreso.
\clearpage\section{Fuentes de apoyo y alcance de consulta}
OpenStax. Introductory Statistics 2e: Chapter 1, Introduction. Consulta: 10 de octubre de 2026. \url{https://openstax.org/books/introductory-statistics-2e/pages/1-introduction}\par Introducción consultada sobre interpretación de datos y muestreo; ejercicios aritméticos y algebraicos originales, verificados por cálculo.\par\medskip
Los originales bibliográficos disponibles se identifican por nombre de archivo y SHA256 en el inventario local. No se atribuyen ediciones, páginas ni lecturas no comprobadas.
\end{document}
